\documentclass[a4paper]{article}

% - - - - - - -  - - - - - Packages - - - - - - - - - - - - 
\usepackage{amsmath}
\usepackage{amsfonts}                       % Add mathematical symbols and functionality
\usepackage{amssymb}
\usepackage[margin=2.2cm]{geometry}         % Change margins
\usepackage[shortlabels]{enumitem}          % Enumerate things
\usepackage{graphicx}                       % Enables pictures 

%\usepackage[english]{babel}
\usepackage[dutch]{babel}                   % Language settings

\usepackage{float}                          % Force pictures in right place
\usepackage{braket}                         % Make braket notation and sets

% - - - - - - -  - - - - - Newcommands - - - - - - - - - - - 
\newcommand{\N}{\mathbb{N}}
\newcommand{\Z}{\mathbb{Z}}
\newcommand{\id}{\mathrm{id}}

\newcommand{\norm}[1]{ \left\| #1 \right\| }

\newcommand{\pointsvector}{\overrightarrow{AB}}


\title{Inlever template}
\author{\TeX niCie}


\begin{document}

\maketitle

Dit is een template voor een inleveropgave. 
Laten we beginnen met opgave 1. 
Dat kan bijvoorbeeld op deze manier: 

\begin{enumerate}[1.]
    \item Dit is opgave 1. Opgave 1 een een aantal deelopgaven. 
    Daarom zullen we nu beginnnen met opgave a.
    \begin{enumerate}[a)]
        \item Dit is opgave a. 
        \item Dit is opgave b.
    \end{enumerate}

    \item In de preamble van dit document zijn een aantal nieuwe 
    commands gedefinieerd, zoals $\N$ en $\Z$. 
    We hebben ook een command gemaakt voor de norm, 
    dit command heeft een argument. 
    Je kunt het op deze manier gebruiken: $\norm{x}$.
    Dit werkt ook voor grotere vergelijkingen:
    \[ \norm{\sum_{n=0}^N \frac{(-1)^nf_n}{n^2 + 7}}_{\infty} \leq 
    \sum_{n=0}^N \norm{\frac{(-1)^nf_n}{n^2 + 7}}_{\infty}.  \]
    Hier zijn een aantal opgaven met commands: 
    \begin{enumerate}[(a)]
        \item Voeg zelf commands toe om $\mathbb{R}$
            en $\mathbb{C}$ korter te schrijven.
        \item Maak een command om $\frac{\mathrm{d}}{\mathrm{d}x}$ 
            korter te schrijven.
        \item In de preamble staat een command $\pointsvector$, dit is nu een vector 
            vanaf $A$ naar $B$. 
            Pas dit aan zodat je zelf de punten kunt geven 
            wanneer je het command aanroept (geef dit command een argument).
    \end{enumerate}

    % Dit is een comment 
    % Deze tekst zal niet in je pdf-document verschijnen

    \item We hebben een aantal packages toegevoegd aan
    aan deze template, bijvoorbeeld de enumitem package, 
    We hebben deze nodig om voor het nummeren van de opgaves. 
    We hebben ook de optie ``shortlabels'' toegevoegd. 
    Hierdoor begrijpt de package dingen als ``[(i)]'' zoals hier:
    \begin{enumerate}[(i)]
        \item We hebben ook een aantal andere packages toegevoegd. 
        \item Voor de natuurkundigen zal de braket package handig worden. 
            Dit voegt notatie zoals $\bra{\phi}$ en $\ket{\psi}$ toe. 
            Voor wiskundigen is deze package ook handig, want het voegt 
            de command $\set{1,2,3}$ toe. 
        \item We hebben ook de package geometry toegevoegd. 
            Dit kun je gebruiken om de marges aan te passen. 
    \end{enumerate}


\end{enumerate}

\end{document}